\section{Background and Preference Orderings}
\label{sec:background}

Before introducing vector-valued relational states, it is necessary to show
precisely why classical preference representations break down under concurrent
functional partitioning. This section reviews standard ordinal structures from
social choice and matching theory and establishes an impossibility result for
multi-role partner portfolios.

\subsection{Total Orders, Weak Orders, and Interval Orders}
\label{subsec:bg-orders}

Let $\partnerset = \{G_1, \dots, G_n\}$ denote a finite set of prospective or
active partners. In classical decision theory~\citep{arrow1963social,fishburn1970utility,sen2017collective},
the focal agent $W$ evaluates alternatives through a binary preference relation
$\succsim$ on $\partnerset$, where $G_i \succsim G_j$ means ``$G_i$ is at least
as preferred as $G_j$.'' Strict preference is defined by $G_i \succ G_j$ iff
$G_i \succsim G_j$ and $\neg(G_j \succsim G_i)$, while indifference is
$G_i \sim G_j$ iff $G_i \succsim G_j$ and $G_j \succsim G_i$.

\begin{definition}[Classical preference hierarchies]
\label{def:order-hierarchy}
A binary relation $\succsim$ on $\partnerset$ is:
\begin{enumerate}[nosep]
  \item a \emph{preorder} if it is reflexive ($G_i \succsim G_i$) and transitive
        ($G_i \succsim G_j$ and $G_j \succsim G_k \implies G_i \succsim G_k$);
  \item a \emph{weak order} if it is a complete preorder (for all $i, j$, either
        $G_i \succsim G_j$ or $G_j \succsim G_i$);
  \item a \emph{strict total order} if it is an antisymmetric weak order,
        inducing a unique bijection $\rho : \partnerset \to \{1, \dots, n\}$
        giving each partner's integer rank.
\end{enumerate}
\end{definition}

Two-sided matching models~\citep{gale1962college,roth1990two} typically assume
strict total orders (or at weakest, weak orders with tie-breaking rules) so that
every partner occupies a definite rank $\rho(G_i) \in \{1, \dots, n\}$. Even when
utility theory relaxes strictness to \emph{semiorders} or \emph{interval orders}
to accommodate just-noticeable-difference thresholds~\citep{fishburn1970utility},
the underlying valuation remains one-dimensional: every partner is mapped to a
real number or interval on a single scalar axis.

\subsection{Ordinal Versus Cardinal Rank Under Multi-Attribute Tradeoffs}
\label{subsec:bg-cardinal}

Suppose instead that $W$ evaluates partners along $m \ge 2$ distinct functional
criteria indexed by roles $\roleset = \{r_1, \dots, r_m\}$ (for instance, public
compatibility, emotional attunement, material support, and weekend leisure).
Each partner $G_i$ realizes a non-negative performance profile
$\mathbf{u}(G_i) = (u_{i1}, \dots, u_{im}) \in \Rnum_{\ge 0}^m$, inducing $m$
separate role-specific weak orders $\succsim_{r_1}, \dots, \succsim_{r_m}$ on
$\partnerset$.

Aggregating $\{\succsim_{r_j}\}_{j=1}^m$ into a single total order $\succsim$
encounters the classic multi-criteria aggregation obstruction analogous to
\citet{arrow1963social}. Specifically, if no single role is dictatorial, and if
$W$ treats partners symmetrically across disjoint role specializations, then
strict Pareto dominance fails to produce a complete ranking: partner $G_1$ leads
on $r_1$, partner $G_2$ leads on $r_2$, and partner $G_3$ leads on $r_3$. None
dominates the others vectorwise, leaving $\succsim$ strictly incomplete.

\begin{proposition}[Failure of scalar rank under functional specialization]
\label{prop:order-failure}
Let $m \ge 2$ and $n \ge 2$. Suppose each partner $G_i \in \partnerset$ attains
strict superiority in at least one distinct role $r_{\phi(i)}$ such that
$u_{i,\phi(i)} > u_{k,\phi(i)}$ for all $k \neq i$. Then:
\begin{enumerate}[nosep]
  \item the Pareto dominance relation $\succsim_{\mathrm{Pareto}}$ has no
        unique maximal element, and all pairs $(G_i, G_k)$ with $i \neq k$ are
        incomparable;
  \item any scalar rank assignment $\rho : \partnerset \to \{1, \dots, n\}$
        compatible with $\succsim_{\mathrm{Pareto}}$ depends on an arbitrary
        role-weighting vector $\boldsymbol{\alpha} \in \Delta^{m-1}$, and
        distinct choices of $\boldsymbol{\alpha}$ permute which partner receives
        rank~$1$.
\end{enumerate}
\end{proposition}

\begin{proof}
For any $i \neq k$, $G_i$ strictly exceeds $G_k$ on coordinate $\phi(i)$, while
$G_k$ strictly exceeds $G_i$ on coordinate $\phi(k)$. Hence neither
$\mathbf{u}(G_i) \ge \mathbf{u}(G_k)$ nor $\mathbf{u}(G_k) \ge \mathbf{u}(G_i)$
holds componentwise, proving (1). For (2), consider linear scalarizations
$S_{\boldsymbol{\alpha}}(G_i) = \sum_{j=1}^m \alpha_j u_{ij}$ with
$\alpha_j \ge 0$ and $\sum_j \alpha_j = 1$. Setting $\boldsymbol{\alpha}$ equal
to the unit basis vector $\mathbf{e}_{\phi(i)}$ makes
$S_{\boldsymbol{\alpha}}(G_i) = u_{i,\phi(i)} > u_{k,\phi(i)} = S_{\boldsymbol{\alpha}}(G_k)$
for all $k \neq i$, placing $G_i$ uniquely at rank~$1$. Since this holds for
every $i \in \{1, \dots, n\}$, no unique rank-1 partner is determined by the
underlying attribute profiles alone.
\end{proof}

\subsection{Why Evasiveness Is Structural Rather Than Strategic}
\label{subsec:bg-evasiveness}

\cref{prop:order-failure} explains a recurring empirical regularity. In
situationships, $W$ does not maintain a secret integer list $(G_{(1)}, G_{(2)}, \dots, G_{(n)})$
that is merely concealed from the partners. Instead, $W$'s internal role weights
$\boldsymbol{\alpha}(t)$ fluctuate across contexts: on Tuesday night during an
interpersonal crisis, weight concentrates on $r_2$ (emotional support), making
$G_2$ effectively primary; on Saturday night, weight concentrates on $r_4$
(entertainment), making $G_4$ effectively primary; when bills or errands arise,
weight shifts toward $r_3$.

Asking ``who is first?'' without specifying a context or forcing a permanent
scalarization $\boldsymbol{\alpha}$ is therefore asking for an invariant that
does not exist prior to measurement. Formalizing this state requires replacing
static integer rankings with linear combinations over a role basis and
probabilistic weights over partners, which we construct in \cref{sec:framework}.
